% Options for packages loaded elsewhere
\PassOptionsToPackage{unicode}{hyperref}
\PassOptionsToPackage{hyphens}{url}
\documentclass[11pt]{article}

\usepackage[margin=1in]{geometry}
\usepackage{xcolor}
\usepackage{amsmath,amssymb}
\usepackage[T1]{fontenc}
\usepackage[utf8]{inputenc}
\usepackage{textcomp}
\usepackage{lmodern}
\usepackage{microtype}
\usepackage{parskip}
\usepackage{color}
\usepackage{fancyvrb}
\usepackage{longtable,booktabs,array}
\usepackage{calc}
\usepackage{etoolbox}
\usepackage{graphicx}
\usepackage{bookmark}
\usepackage{xurl}
\usepackage{hyperref}

% Suppress section auto-numbering (headings carry manual numbers)
\setcounter{secnumdepth}{-\maxdimen}

% Code highlighting (generated by Pandoc)
\newcommand{\VerbBar}{|}
\newcommand{\VERB}{\Verb[commandchars=\\\{\}]}
\DefineVerbatimEnvironment{Highlighting}{Verbatim}{commandchars=\\\{\},fontsize=\small}
\newenvironment{Shaded}{\small}{}
\newcommand{\AlertTok}[1]{\textcolor[rgb]{1.00,0.00,0.00}{\textbf{#1}}}
\newcommand{\AnnotationTok}[1]{\textcolor[rgb]{0.38,0.63,0.69}{\textbf{\textit{#1}}}}
\newcommand{\AttributeTok}[1]{\textcolor[rgb]{0.49,0.56,0.16}{#1}}
\newcommand{\BaseNTok}[1]{\textcolor[rgb]{0.25,0.63,0.44}{#1}}
\newcommand{\BuiltInTok}[1]{\textcolor[rgb]{0.00,0.50,0.00}{#1}}
\newcommand{\CharTok}[1]{\textcolor[rgb]{0.25,0.44,0.63}{#1}}
\newcommand{\CommentTok}[1]{\textcolor[rgb]{0.38,0.63,0.69}{\textit{#1}}}
\newcommand{\CommentVarTok}[1]{\textcolor[rgb]{0.38,0.63,0.69}{\textbf{\textit{#1}}}}
\newcommand{\ConstantTok}[1]{\textcolor[rgb]{0.53,0.00,0.00}{#1}}
\newcommand{\ControlFlowTok}[1]{\textcolor[rgb]{0.00,0.44,0.13}{\textbf{#1}}}
\newcommand{\DataTypeTok}[1]{\textcolor[rgb]{0.56,0.13,0.00}{#1}}
\newcommand{\DecValTok}[1]{\textcolor[rgb]{0.25,0.63,0.44}{#1}}
\newcommand{\DocumentationTok}[1]{\textcolor[rgb]{0.73,0.13,0.13}{\textit{#1}}}
\newcommand{\ErrorTok}[1]{\textcolor[rgb]{1.00,0.00,0.00}{\textbf{#1}}}
\newcommand{\ExtensionTok}[1]{#1}
\newcommand{\FloatTok}[1]{\textcolor[rgb]{0.25,0.63,0.44}{#1}}
\newcommand{\FunctionTok}[1]{\textcolor[rgb]{0.02,0.16,0.49}{#1}}
\newcommand{\ImportTok}[1]{\textcolor[rgb]{0.00,0.50,0.00}{\textbf{#1}}}
\newcommand{\InformationTok}[1]{\textcolor[rgb]{0.38,0.63,0.69}{\textbf{\textit{#1}}}}
\newcommand{\KeywordTok}[1]{\textcolor[rgb]{0.00,0.44,0.13}{\textbf{#1}}}
\newcommand{\NormalTok}[1]{#1}
\newcommand{\OperatorTok}[1]{\textcolor[rgb]{0.40,0.40,0.40}{#1}}
\newcommand{\OtherTok}[1]{\textcolor[rgb]{0.00,0.44,0.13}{#1}}
\newcommand{\PreprocessorTok}[1]{\textcolor[rgb]{0.74,0.48,0.00}{#1}}
\newcommand{\RegionMarkerTok}[1]{#1}
\newcommand{\SpecialCharTok}[1]{\textcolor[rgb]{0.25,0.44,0.63}{#1}}
\newcommand{\SpecialStringTok}[1]{\textcolor[rgb]{0.73,0.40,0.53}{#1}}
\newcommand{\StringTok}[1]{\textcolor[rgb]{0.25,0.44,0.63}{#1}}
\newcommand{\VariableTok}[1]{\textcolor[rgb]{0.10,0.09,0.49}{#1}}
\newcommand{\VerbatimStringTok}[1]{\textcolor[rgb]{0.25,0.44,0.63}{#1}}
\newcommand{\WarningTok}[1]{\textcolor[rgb]{0.38,0.63,0.69}{\textbf{\textit{#1}}}}

% Table fixes
\makeatletter
\patchcmd\longtable{\par}{\if@noskipsec\mbox{}\fi\par}{}{}
\makeatother
\IfFileExists{footnotehyper.sty}{\usepackage{footnotehyper}}{\usepackage{footnote}}
\makesavenoteenv{longtable}

% Figure scaling
\makeatletter
\newsavebox\pandoc@box
\newcommand*\pandocbounded[1]{%
  \sbox\pandoc@box{#1}%
  \Gscale@div\@tempa{\textheight}{\dimexpr\ht\pandoc@box+\dp\pandoc@box\relax}%
  \Gscale@div\@tempb{\linewidth}{\wd\pandoc@box}%
  \ifdim\@tempb\p@<\@tempa\p@\let\@tempa\@tempb\fi%
  \ifdim\@tempa\p@<\p@\scalebox{\@tempa}{\usebox\pandoc@box}%
  \else\usebox{\pandoc@box}%
  \fi%
}
\def\fps@figure{htbp}
\makeatother

\setlength{\emergencystretch}{3em}
\providecommand{\tightlist}{%
  \setlength{\itemsep}{0pt}\setlength{\parskip}{0pt}}

\hypersetup{hidelinks, pdfcreator={LaTeX}}

% ---------------------------------------------------------------
\title{\textbf{Distributed KV Cache Retrieval for RAG via Kademlia P2P Networks:
Simulation-Validated Optimal Splitting Strategies}}
\author{Joseph Huynh\\
  California State University, Long Beach\\
  \texttt{CECS497 - Directed Studies}}
\date{May 2026}
% ---------------------------------------------------------------

\begin{document}

\maketitle

\begin{abstract}
Time-to-first-token (TTFT) in retrieval-augmented generation (RAG) is
bottlenecked by loading precomputed KV caches sequentially from a single
host. This paper investigates whether distributing KV cache chunks
across a peer-to-peer (P2P) Kademlia network and fetching layer-group
stripes in parallel can reduce this bottleneck. We profile KV cache
memory for GPT-2 (73,728 bytes/token, 12 layers) and TinyLLaMA-1.1B
(45,056 bytes/token, 22 layers, GQA), design a binary serialization
format with optional compression, and select Kademlia as the DHT
backbone based on its 100\% lookup correctness and $\alpha$-parallel iterative
lookup. A three-stage simulation experiment (120 + 4,050 + 3,240 = 7,410
total runs) evaluates 12 chunk-splitting strategies across network sizes
of 32--128 peers, replication factors of 1--3, churn rates of 0--15\%,
and context lengths of 128--1,024 tokens. The optimal configuration ---
\texttt{token\_block=32}, \texttt{layer\_group=2} --- achieves a confirmatory p95
retrieval latency of \textbf{4.911\,ms $\pm$ 0.032\,ms} at \textbf{99.56\%
chunk availability}, a 26\% improvement over the next-best variant and
3--7$\times$ improvement over one-dimensional splitting strategies. All results
are stable across the full parameter sweep. A single-machine P2P
prototype with five implementation milestones is fully specified and
ready for construction; its success criterion is a statistically
significant p95 TTFT reduction over a sequential-load baseline across
200 queries.
\end{abstract}

% ---------------------------------------------------------------
\section{1. Introduction}

Retrieval-augmented generation augments LLM inference by prepending
relevant document context to each query. In systems that cache the KV
representations of documents offline (e.g., RAGCache, CacheBlend), the
TTFT for a RAG query consists of: (1) retrieval of relevant document
IDs, and (2) loading the precomputed KV cache for those documents into
GPU memory before decoding. On a single host, step (2) is a sequential
disk or memory read bounded by local I/O bandwidth.

The central hypothesis of this work is: \emph{if a document's KV cache
is split into stripes held by different peers, fetching all stripes in
parallel exploits aggregate network bandwidth and reduces the
critical-path latency compared to a single sequential load.} The design
challenge is identifying the right splitting granularity --- too coarse
and parallelism is lost; too fine and DHT routing overhead dominates.

\textbf{Contributions:}

\begin{enumerate}
\tightlist
\item
  \textbf{KV cache size profile:} TinyLLaMA-1.1B has 45,056 bytes/token
  (22.1\,MB at a 512-token context); GPT-2 has 73,728 bytes/token (36.0\,MB
  at 512 tokens). An 8\,GiB budget accommodates $\sim$190,650
  TinyLLaMA tokens in KV cache.
\item
  \textbf{Binary serialization format:} A 64-byte chunk header with
  optional lz4/zstd compression. Real TinyLLaMA KV data compresses 5.8\%
  with zstd at a 2.26$\times$ roundtrip time cost; lz4 offers negligible
  compression at 1.21$\times$ cost.
\item
  \textbf{DHT selection:} Kademlia outperforms Chord (100\% vs.\ 100\%
  correctness, lower join cost) and dominates Pastry (100\% vs.\ 91.5\%
  correctness). Its $\alpha$-parallel iterative lookup maps directly to
  multi-stripe fetching.
\item
  \textbf{Optimal splitting strategy:} A three-stage simulation
  experiment (7,410 runs) confirms
  \texttt{token\_block=32}, \texttt{layer\_group=2} as Pareto-dominant: p95 =
  4.911\,ms $\pm$ 0.032\,ms, 99.56\% availability, robust across all tested
  conditions.
\item
  \textbf{P2P system design:} A five-milestone single-machine prototype
  specification with defined acceptance tests and a success criterion.
\end{enumerate}

% ---------------------------------------------------------------
\section{2. Background}

\subsection{2.1 KV Caching and RAG}

Autoregressive LLMs compute key and value tensors for every token in the
context. These tensors, collectively the KV cache, grow linearly with
sequence length and model depth. In RAG, document context is fixed per
document and can be prefilled offline; at query time, only the user's
question needs new computation if the document KV cache is loaded
intact. RAGCache~[4] and CacheBlend~[5] exploit this property on a
single host. This work extends the paradigm to a P2P setting where each
peer holds a layer-group stripe of a document's KV cache.

\subsection{2.2 Related Systems}

\textbf{RAGCache}~[4] stores full KV caches on a single node; retrieval
is sequential. \textbf{CacheBlend}~[5] merges KV caches from multiple
documents with a blending step to correct for positional interference.
\textbf{TurboRAG}~[6] encodes document chunks with a fixed dummy-position
offset (Reordered-RoPE), making chunks position-invariant and eliminating
the need for a stitching pass at retrieval time. Our design uses standard
RoPE and requires a stitching correction at retrieval; the TurboRAG
approach is noted as a fallback if the correction proves numerically
unstable.

\subsection{2.3 Kademlia DHT}

Kademlia~[1] routes lookups using XOR distance on a 160-bit key space.
Each node maintains a k-bucket routing table partitioned by XOR distance.
Lookups are \emph{iterative} and \emph{parallel}: $\alpha$ probes are
issued simultaneously at each round, converging in $O(\log N)$ hops. The
$\alpha$-parallel structure means that fetching $k$ stripes from $k$ peers
can be pipelined within a single lookup round, which is the key property
exploited by the swarm fetcher in Section~3.3.

\subsection{2.4 Rotary Position Embedding (RoPE)}

TinyLLaMA uses RoPE, which encodes absolute positions into key and query
tensors via rotation matrices. When a document chunk is prefilled at
positions $[0, L_\text{doc})$, its keys carry those absolute rotations.
To stitch the chunk into a query context at positions $[L_\text{query},
L_\text{query} + L_\text{doc})$, the keys must be reverse-rotated and
re-rotated at the new positions. This stitching pass is the
correctness-critical step in the P2P prototype (Section~8).

% ---------------------------------------------------------------
\section{3. Design}

\subsection{3.1 DHT Protocol Selection}

Three DHT protocols were evaluated in a virtual-time simulation (N=20
nodes, 200 random key lookups per protocol). The selection criterion
prioritized lookup correctness first, then latency and join cost.

\begin{longtable}[]{@{}lllllll@{}}
\toprule
Protocol & Mean hops & P95 hops & Max hops & Correctness & Join msgs & Routing entries \\
\midrule
\endhead
\bottomrule
\endlastfoot
Chord & 1.90 & 4 & 4 & \textbf{100\%} & 16.0 & 4.8 \\
\textbf{Kademlia} & \textbf{1.61} & \textbf{3} & \textbf{3} & \textbf{100\%} & \textbf{2.4} & 12.5 \\
Pastry & 1.12 & 2 & 2 & 91.5\% & 7.0 & 19.9 \\
\end{longtable}

Pastry~[3] achieves the lowest hop count but fails 8.5\% of lookups at
N=20 and degrades to 9\% failure post-churn --- unacceptable for a
system where every missing stripe blocks decode. Chord~[2] and
Kademlia~[1] both achieve 100\% correctness; Kademlia is selected because
its $\alpha$-parallel iterative lookup ($\alpha$=3 in experiments) issues
multiple probes simultaneously, a property that maps directly to
concurrent stripe fetching. Chord's sequential ring traversal does not
offer this.

\subsection{3.2 KV Cache Serialization Format}

Each chunk is serialized as a 64-byte header followed by a raw payload:

\begin{longtable}[]{@{}lll@{}}
\toprule
Field & Size & Description \\
\midrule
\endhead
\bottomrule
\endlastfoot
Magic bytes & 4\,B & \texttt{KVC1} --- format identifier \\
Version & 1\,B & Format version \\
Compression codec & 1\,B & 0=none, 1=lz4, 2=zstd \\
Data type & 1\,B & torch dtype enum \\
Token start / end & 8\,B & Inclusive range \\
Layer start / end & 8\,B & Inclusive range \\
Tensor shape & 16\,B & (layers, heads, tokens, head\_dim) \\
MD5 checksum & 16\,B & Checksum of uncompressed payload \\
Reserved & 9\,B & Future use \\
\end{longtable}

The chunk ID encodes the provenance:
\texttt{m=<model>|ts=<tok\_start>|te=<tok\_end>|ls=<lay\_start>|le=<lay\_end>}.
Default splitting parameters --- \texttt{token\_block=32},
\texttt{layer\_group=2} --- were determined by the simulation experiment
in Section~6.3 and are applied in \texttt{kv\_serialization/serialize.py}.

Compression was evaluated empirically on real TinyLLaMA KV data
(Section~6.2). lz4 is available as an optional fast path; zstd is
available for bandwidth-constrained deployments.

\subsection{3.3 P2P System Architecture}

The full prototype consists of five components:

\begin{verbatim}
Query Node
  +-- RAG Retriever  --> top-k doc IDs
  +-- Swarm Fetcher  --> Peer 0: layers  0-5   (stripe 0)
                    --> Peer 1: layers  6-11  (stripe 1) --> RoPE Stitcher --> LLM Decode
                    --> Peer 2: layers 12-17  (stripe 2)
                    --> Peer 3: layers 18-21  (stripe 3)
\end{verbatim}

Each peer serves a single layer-group stripe for each cached document
via a lightweight HTTP endpoint
(\texttt{GET /chunk?doc\_id=X\&layer\_group=Y}). The swarm fetcher
issues all four requests concurrently via \texttt{asyncio}/\texttt{aiohttp}
and reassembles stripes in layer order. The RoPE stitcher then corrects
position encodings before handing the assembled KV cache to the LLM
decoder, skipping the full prefill pass.

The speculative prefetcher (M4) pre-fetches predicted top-k document
stripes into a 16-entry LRU cache while a prior request is being
decoded, hiding fetch latency under inference time.

Planned deliverables: \texttt{scripts/peer\_server.py}~(M1),
\texttt{kv\_kademlia\_experiments/swarm\_fetcher.py}~(M2),
\texttt{kv\_kademlia\_experiments/rope\_stitcher.py}~(M3),
\texttt{kv\_kademlia\_experiments/speculative\_prefetcher.py}~(M4),
\texttt{scripts/benchmark\_e2e\_ttft.py}~(M5).

% ---------------------------------------------------------------
\section{4. Environment Setup and Running the Code}

\subsection{4.1 Prerequisites}

\begin{itemize}
\tightlist
\item Python 3.10 or later. CUDA is optional; all scripts fall back to CPU.
\item Create and activate a virtual environment:
\end{itemize}

\begin{Shaded}
\begin{Highlighting}[]
\CommentTok{\# Windows}
\ExtensionTok{python} \AttributeTok{-m} venv venv
\ExtensionTok{venv\textbackslash{}Scripts\textbackslash{}activate}

\CommentTok{\# Unix / macOS}
\ExtensionTok{python} \AttributeTok{-m} venv venv
\BuiltInTok{source} venv/bin/activate
\end{Highlighting}
\end{Shaded}

\begin{itemize}
\tightlist
\item Install dependencies:
\end{itemize}

\begin{Shaded}
\begin{Highlighting}[]
\ExtensionTok{pip} install \AttributeTok{-r} requirements.txt
\end{Highlighting}
\end{Shaded}

The planned P2P system requires two additional packages not yet in
\texttt{requirements.txt}:

\begin{Shaded}
\begin{Highlighting}[]
\ExtensionTok{pip} install aiohttp fastapi uvicorn
\end{Highlighting}
\end{Shaded}

\subsection{4.2 Running Each Script}

\begin{longtable}[]{@{}llll@{}}
\toprule
Script & Command & Expected runtime & Output location \\
\midrule
\endhead
\bottomrule
\endlastfoot
GPT-2 KV profiling & \texttt{python scripts/kv\_cache\_phase1\_gpt2.py} & $\sim$2\,min (CPU) & \texttt{phase1\_outputs/} \\
TinyLLaMA KV profiling & \texttt{python scripts/kv\_cache\_phase2\_tinyllama.py} & $\sim$5\,min (CPU) & \texttt{phase2\_outputs/} \\
Real KV compression & \texttt{python scripts/benchmark\_real\_kv\_compression.py} & $\sim$3\,min & \texttt{benchmark\_outputs/} \\
Serialization benchmark & \texttt{python scripts/benchmark\_serialization.py} & $<$1\,min & \texttt{benchmark\_outputs/} \\
Kademlia experiment & \texttt{python scripts/experiment\_tinyllama\_kv\_kademlia.py} & A $\sim$20\,min; B $\sim$4\,h; C $\sim$8\,h & \texttt{kv\_kademlia\_experiments/} \\
Visualization & \texttt{python scripts/viz\_kademlia\_latency.py} & $<$1\,min & \texttt{figures/} \\
\end{longtable}

The Kademlia experiment script runs all three stages sequentially by
default. Individual stages can be run by passing \texttt{--stage a},
\texttt{--stage b}, or \texttt{--stage c}. Stages B and C read the
winner from the prior stage's \texttt{variant\_ranking.csv}
automatically.

\subsection{4.3 Planned P2P System Commands}

Once M1--M5 are implemented, the full system runs as follows:

\begin{Shaded}
\begin{Highlighting}[]
\CommentTok{\# Start four peer processes (one per layer-group stripe)}
\ExtensionTok{python} scripts/peer\_server.py \AttributeTok{--port} 5001 \AttributeTok{--layer-group} 0
\ExtensionTok{python} scripts/peer\_server.py \AttributeTok{--port} 5002 \AttributeTok{--layer-group} 1
\ExtensionTok{python} scripts/peer\_server.py \AttributeTok{--port} 5003 \AttributeTok{--layer-group} 2
\ExtensionTok{python} scripts/peer\_server.py \AttributeTok{--port} 5004 \AttributeTok{--layer-group} 3

\CommentTok{\# Run the end-to-end TTFT benchmark}
\ExtensionTok{python} scripts/benchmark\_e2e\_ttft.py
\end{Highlighting}
\end{Shaded}

Output: \texttt{P2P\_RAGCACHE\_RESULTS.md} and a companion CSV with
per-query TTFT measurements.

% ---------------------------------------------------------------
\section{5. Testing}

\subsection{5.1 Running the Test Suite}

\begin{Shaded}
\begin{Highlighting}[]
\ExtensionTok{pytest} tests/ \AttributeTok{-v} \AttributeTok{-s}
\end{Highlighting}
\end{Shaded}

All 54 tests complete in approximately 2 seconds.

\begin{longtable}[]{@{}lll@{}}
\toprule
Test file & Coverage & Tests \\
\midrule
\endhead
\bottomrule
\endlastfoot
\texttt{tests/test\_protocols.py} & Chord, Kademlia, Pastry correctness; join/leave; churn & 27 \\
\texttt{tests/test\_serialization.py} & Round-trip; lz4/zstd; MD5 integrity; reassembly & 21 \\
\texttt{tests/benchmark\_kademlia\_latency.py} & $O(\log N)$ hop-count scaling & 6 \\
\end{longtable}

\subsection{5.2 Planned Acceptance Tests for P2P Components}

Each milestone has a defined pass/fail criterion:

\begin{itemize}
\tightlist
\item
  \textbf{M1 (peer server):} A \texttt{requests.get('http://localhost:5001/chunk?doc\_id=0\&layer\_group=0')}
  returns the correct bytes and round-trips to a valid tensor with matching MD5.
\item
  \textbf{M2 (swarm fetcher):} p95 fetch latency over 100 requests $<$ 50\,ms on localhost.
\item
  \textbf{M3 (RoPE stitcher):} Cosine similarity between stitched-KV forward-pass logits
  and a reference full-prefill forward pass $>$ 0.99 on the final hidden state (fp16 tolerance).
\item
  \textbf{M5 (TTFT benchmark):} p95 TTFT lower than sequential-load baseline at
  $p < 0.05$ (Mann--Whitney U) across $\geq 200$ queries.
\end{itemize}

% ---------------------------------------------------------------
\section{6. Results}

\subsection{6.1 KV Cache Size Profile}

KV cache memory scales linearly with sequence length in both models.
TinyLLaMA uses grouped-query attention (GQA) with 4 KV heads instead of
32 full heads, which is the primary reason its per-token cost is 61.1\%
of GPT-2 despite having nearly twice the layers.

\begin{longtable}[]{@{}lllllll@{}}
\toprule
Model & Layers & KV heads & Bytes/token & @ 128 tokens & @ 512 tokens & @ 1,024 tokens \\
\midrule
\endhead
\bottomrule
\endlastfoot
GPT-2 & 12 & 12 & 73,728 & 9.44\,MB & 36.0\,MB & 72.0\,MB \\
TinyLLaMA-1.1B & 22 & 4 (GQA) & 45,056 & 5.77\,MB & 22.1\,MB & 44.1\,MB \\
Ratio (LLaMA/GPT-2) & --- & --- & \textbf{0.611} & 0.611 & 0.611 & 0.611 \\
\end{longtable}

At 512 tokens --- the representative RAG document chunk size used in all
subsequent experiments --- one TinyLLaMA document's KV cache occupies
\textbf{22.1\,MB uncompressed}. An 8\,GiB GPU memory budget accommodates
approximately 190,650 tokens ($\approx$373 document chunks of 512 tokens each)
in KV cache alone.

\subsection{6.2 Compression}

Compression was evaluated on actual TinyLLaMA KV tensors extracted from
real inference runs. KV cache data is largely random-looking
(floating-point activations), which limits compression effectiveness.

\begin{longtable}[]{@{}llllll@{}}
\toprule
Sequence length & Codec & Raw bytes & Compressed bytes & Ratio & Roundtrip time \\
\midrule
\endhead
\bottomrule
\endlastfoot
512 tokens & none  & 23,068,672 & 23,079,936 & 1.000$\times$ & 83.0\,ms \\
512 tokens & lz4   & 23,068,672 & 23,061,956 & 1.001$\times$ & 100.2\,ms \\
512 tokens & zstd  & 23,068,672 & 21,743,156 & \textbf{1.061$\times$} & 187.4\,ms \\
\end{longtable}

\textbf{Key finding:} zstd achieves a consistent \textbf{1.061$\times$
compression ratio} (5.8\% size reduction) across all tested sequence
lengths (32--512 tokens), at the cost of a \textbf{2.26$\times$ roundtrip
time overhead} (187\,ms vs.\ 83\,ms at 512 tokens). lz4 offers effectively
zero compression on this data type. Neither codec is recommended by
default; zstd is a viable option in bandwidth-constrained deployments
where the size reduction outweighs the latency penalty.

\subsection{6.3 Kademlia KV Splitting: Three-Stage Simulation}

The core experiment evaluates 12 chunk-splitting strategies by simulating
KV cache placement and retrieval on a Kademlia network. Simulation
parameters: 16-bit ID space, k-bucket size 8, $\alpha$=3, per-hop delay
1.0\,ms, base transfer time 0.2\,ms, bandwidth 100\,Mbps. Strategies are
scored as: $0.60 \times (1 - \text{norm.\ p95 latency}) + 0.25 \times
\text{availability} + 0.15 \times (1 - \text{overhead ratio})$.

\subsubsection{Stage A: Screening (120 runs)}

Twelve strategy variants were screened at N=64, R=2, churn
$\in \{0\%,\,5\%\}$, L=512, 5 random seeds. One-dimensional strategies
were immediately eliminated:

\begin{longtable}[]{@{}llll@{}}
\toprule
Strategy & p95 latency & Chunk success & Weighted score \\
\midrule
\endhead
\bottomrule
\endlastfoot
\textbf{token\_layer\_group (winner)} & \textbf{8.43\,ms} & \textbf{99.92\%} & \textbf{0.750} \\
layer\_only & 27.44\,ms & 100.0\% & 0.442 \\
token\_only & 36.36\,ms & 100.0\% & 0.250 \\
\end{longtable}

Token-only splitting produces the largest chunks (fewer but larger
stripes), saturating the 100\,Mbps transfer model. Layer-only splitting
produces many small stripes that each require independent routing,
multiplying DHT overhead. The 2D token+layer-group family advanced to
Stage B with four hyperparameter variants.

\subsubsection{Stage B: Full Parameter Sweep (4,050 runs)}

The top five variants from Stage A were evaluated across N $\in
\{32,\,64,\,128\}$, R $\in \{1,\,2,\,3\}$, churn $\in \{0\%,\,5\%,\,15\%\}$,
L $\in \{128,\,512,\,1024\}$, 10 seeds.

\begin{longtable}[]{@{}llll@{}}
\toprule
Variant & p95 latency & Availability & Weighted score \\
\midrule
\endhead
\bottomrule
\endlastfoot
\textbf{tb=32, lg=2} & \textbf{4.92\,ms} & \textbf{99.66\%} & \textbf{0.850} \\
tb=32, lg=4 & 6.20\,ms & 99.28\% & 0.734 \\
tb=64, lg=2 & 6.28\,ms & 99.15\% & 0.712 \\
tb=64, lg=4 & 8.82\,ms & 99.11\% & 0.551 \\
layer\_only, lb=1 & 14.84\,ms & 98.75\% & 0.150 \\
\end{longtable}

The \texttt{tb=32, lg=2} variant wins on both latency and availability.
Smaller token blocks (32 vs.\ 64) reduce per-stripe transfer time; finer
layer groups (lg=2 = 11 parallel stripes for TinyLLaMA's 22 layers)
maximize the critical-path reduction from parallel fetching.

\subsubsection{Stage C: Confirmatory (3,240 runs)}

The top two variants were re-evaluated at 20 seeds to obtain tight
confidence intervals.

\begin{longtable}[]{@{}llllll@{}}
\toprule
Variant & p50 & p95 & p99 & Availability & 95\% CI on p95 \\
\midrule
\endhead
\bottomrule
\endlastfoot
\textbf{tb=32, lg=2} & \textbf{3.779\,ms} & \textbf{4.911\,ms} & \textbf{5.293\,ms} & \textbf{99.56\%} & \textbf{$\pm$0.032\,ms} \\
tb=32, lg=4 & 4.994\,ms & 6.190\,ms & 6.558\,ms & 99.19\% & $\pm$0.032\,ms \\
\end{longtable}

The \texttt{tb=32, lg=2} configuration is confirmed as Pareto-dominant:
\textbf{26\% lower p95 latency} than the runner-up, \textbf{37 basis
points higher availability}, and non-overlapping confidence intervals.

\subsubsection{Sensitivity Analysis}

The winner's p95 latency was stable across all parameter dimensions
tested in Stage B and confirmed in Stage C:

\begin{longtable}[]{@{}llll@{}}
\toprule
Parameter & Range & p95 range (tb=32, lg=2) & Interpretation \\
\midrule
\endhead
\bottomrule
\endlastfoot
Network size N & 32 $\to$ 128 & 4.27 $\to$ 5.54\,ms & Sub-linear; consistent with $O(\log N)$ routing \\
Replication R & 1 $\to$ 3    & 5.02 $\to$ 4.82\,ms & Modest $\sim$4\% improvement; diminishing past R=2 \\
Churn & 0\% $\to$ 15\%       & 5.01 $\to$ 4.81\,ms & Stable; spare replication absorbs node loss \\
Sequence length L & 128 $\to$ 1,024 & 5.15 $\to$ 4.72\,ms & Slight improvement; larger chunks amortize routing \\
\end{longtable}

The stability across churn rates (latency \emph{improves} slightly at
higher churn) is a simulation artifact: with R=2, spare replicas are
always available, and Kademlia routing naturally selects the closer
replica as peers leave. This will be re-evaluated in the real-system
prototype.

\begin{figure}[htbp]
\centering
\pandocbounded{\includegraphics[width=\linewidth]{results/kv_kademlia/stage_c_final/figures/latency_p95_by_variant.png}}
\caption{p95 retrieval latency by variant (Stage C, 3,240 runs). \texttt{tb=32 lg=2} has non-overlapping confidence intervals with all other variants.}
\end{figure}

\begin{figure}[htbp]
\centering
\pandocbounded{\includegraphics[width=\linewidth]{results/kv_kademlia/stage_c_final/figures/latency_vs_num_peers_by_variant.png}}
\caption{p95 latency vs.\ network size N for the top two variants. Sub-linear growth is consistent with $O(\log N)$ Kademlia routing.}
\end{figure}

\begin{figure}[htbp]
\centering
\pandocbounded{\includegraphics[width=\linewidth]{results/kv_kademlia/stage_c_final/figures/availability_vs_churn.png}}
\caption{Chunk availability vs.\ churn rate by replication factor. Both variants remain above 99\% availability across all tested churn conditions.}
\end{figure}

\begin{figure}[htbp]
\centering
\pandocbounded{\includegraphics[width=\linewidth]{results/figures/kademlia_latency_scaling.png}}
\caption{Mean Kademlia lookup latency vs.\ network size with $O(\log N)$ reference curve, confirming expected DHT scaling behavior.}
\end{figure}

% ---------------------------------------------------------------
\section{7. Discussion}

The simulation results establish a clear hierarchy among splitting
strategies. The 2D (token + layer-group) family dominates 1D approaches
by 3--7$\times$ in p95 latency because it creates stripes that are
simultaneously (a) small enough to transfer quickly and (b) numerous
enough to fill the $\alpha$-parallel Kademlia routing pipeline. Within the
2D family, the performance difference between \texttt{lg=2} and
\texttt{lg=4} (4.91 vs.\ 6.19\,ms) suggests that finer layer grouping
increases DHT routing overhead faster than it reduces per-stripe transfer
time, so there is an optimal granularity that the simulated 100\,Mbps
model places at \texttt{lg=2}.

One caveat applies to the churn stability finding: the simulation assigns
fixed spare replicas and assumes instantaneous failover routing. A real
system with correlated failures or cold-cache peers may show different
behavior. The speculative prefetcher (M4) is designed to partially
mitigate this by pre-loading stripes during the prior decode window.

% ---------------------------------------------------------------
\section{8. Future Work: P2P RAGCache Prototype}

\subsection{8.1 Gap Between Simulation and Publication}

The simulation establishes that \texttt{tb=32, lg=2} is the right
splitting strategy. What remains is a real TTFT measurement: does the
latency reduction in simulation (sub-5\,ms retrieval) translate to a
measurable TTFT improvement when the full pipeline --- network I/O, RoPE
stitching, and LLM decode --- is exercised end-to-end?

\subsection{8.2 Implementation Milestones}

\begin{longtable}[]{@{}llll@{}}
\toprule
Milestone & Deliverable & Depends on & Duration \\
\midrule
\endhead
\bottomrule
\endlastfoot
M1: Peer server & \texttt{scripts/peer\_server.py} & --- & 1--2 days \\
M2: Swarm fetcher & \texttt{kv\_kademlia\_experiments/swarm\_fetcher.py} & M1 & 1 day \\
M3: RoPE stitcher & \texttt{kv\_kademlia\_experiments/rope\_stitcher.py} & M2 & 2--3 days \\
M4: Speculative prefetcher & \texttt{kv\_kademlia\_experiments/speculative\_prefetcher.py} & M2 & 1 day (parallel to M3) \\
M5: TTFT benchmark & \texttt{scripts/benchmark\_e2e\_ttft.py} & M3, M4 & 2 days \\
\end{longtable}

M1 $\to$ M2 $\to$ M3 $\to$ M5 must be sequential. M4 can be built in
parallel with M3 and merged before M5.

\subsection{8.3 Benchmark Protocol}

\begin{enumerate}
\tightlist
\item Pre-cache KV tensors for 50 document chunks (TinyLLaMA forward passes
  on representative text, serialized with the existing pipeline).
\item Distribute shards to four peer processes (one per layer-group stripe).
\item Issue 200 RAG queries sequentially; each retrieves 1--3 documents.
\item Record \texttt{ttft\_baseline\_ms} (sequential single-file load) and
  \texttt{ttft\_p2p\_ms} (swarm fetch + stitch) per query.
\item Report p50, p95, p99, and mean for both conditions; test significance
  with Mann--Whitney U ($\alpha$=0.05).
\end{enumerate}

\textbf{Success criterion:} p95 TTFT is statistically significantly
lower in the P2P condition than the sequential baseline across
$\geq$200 queries.

\subsection{8.4 Primary Risk}

The RoPE stitching inverse (reverse-rotating key tensors and re-applying
rotations at shifted positions) may accumulate floating-point error. If
cosine similarity between stitched and reference logits falls below 0.99
in the M3 acceptance test, the fallback is to adopt TurboRAG's
fixed-offset encoding: document chunks are prefilled with positions
$[\textit{offset},\, \textit{offset} + L_\text{doc})$ instead of
$[0, L_\text{doc})$, where \textit{offset} is a fixed large constant.
This makes chunks position-invariant at storage time and eliminates the
need for an inverse transform at retrieval time.

% ---------------------------------------------------------------
\section{9. Conclusion}

This paper demonstrates that two-dimensional KV cache splitting --- by
token block and layer group simultaneously --- outperforms
one-dimensional alternatives by 3--7$\times$ in P2P retrieval latency.
The specific configuration \texttt{token\_block=32, layer\_group=2}
achieves a confirmed p95 latency of 4.911\,ms $\pm$ 0.032\,ms at 99.56\%
chunk availability across 7,410 simulation runs spanning the full
parameter space of practical deployment conditions. The result is robust:
it holds across network sizes of 32--128 peers, replication factors of
1--3, churn rates up to 15\%, and document context lengths up to 1,024
tokens.

Kademlia is the right DHT backbone for this system: its 100\% lookup
correctness and $\alpha$-parallel iterative lookup directly enable the
multi-stripe concurrent fetching that produces the latency reduction.
Among all evaluated protocols, it is the only one that combines
correctness guarantees with an architecture that maps naturally to
parallel data retrieval.

The remaining open question --- whether this simulation result translates
to a measurable TTFT improvement in a real system --- is the subject of
the specified P2P RAGCache prototype (Section~8). If the real TTFT
improvement matches the simulation's retrieval latency improvement, the
system contribution would be a concrete, empirically validated approach
to reducing inference latency in retrieval-augmented LLM deployments.

% ---------------------------------------------------------------
\begin{thebibliography}{9}

\bibitem{kademlia}
Petar Maymounkov and David Mazi\`{e}res. 2002.
Kademlia: A peer-to-peer information system based on the XOR metric.
In \textit{Proc.\ 1st Int'l Workshop on Peer-to-Peer Systems (IPTPS '02)}.

\bibitem{chord}
Ion Stoica, Robert Morris, David Karger, M.\ Frans Kaashoek, and Hari Balakrishnan. 2001.
Chord: A scalable peer-to-peer lookup service for internet applications.
In \textit{Proc.\ ACM SIGCOMM '01}, 149--160.

\bibitem{pastry}
Antony Rowstron and Peter Druschel. 2001.
Pastry: Scalable, decentralized object location, and routing for large-scale peer-to-peer systems.
In \textit{Proc.\ IFIP/ACM Middleware '01}, 329--350.

\bibitem{ragcache}
Yuxin Jin et al. 2024.
RAGCache: Efficient Knowledge Caching for Retrieval-Augmented Generation.
\textit{arXiv:2404.12457}.

\bibitem{cacheblend}
Jiayi Yao et al. 2024.
CacheBlend: Fast Large Language Model Serving for RAG with Cached Knowledge Fusion.
\textit{arXiv:2405.16444}.

\bibitem{turborag}
Bin Lu et al. 2024.
TurboRAG: Accelerating Retrieval-Augmented Generation with Precomputed KV Caches for Chunked Text.
\textit{arXiv:2410.23343}.

\bibitem{tinyllama}
Peiyuan Zhang et al. 2024.
TinyLlama: An Open-Source Small Language Model.
\textit{arXiv:2401.02385}.

\bibitem{vllm}
Woojae Kwon et al. 2023.
Efficient Memory Management for Large Language Model Serving with PagedAttention.
In \textit{Proc.\ ACM SOSP '23}.

\bibitem{huggingface}
Thomas Wolf et al. 2020.
HuggingFace Transformers: State-of-the-art Natural Language Processing.
In \textit{Proc.\ EMNLP '20}.

\end{thebibliography}

% ---------------------------------------------------------------
\appendix
\section{A. Experiment Parameter Grids}

\textbf{Stage A} (120 runs):

\begin{longtable}[]{@{}ll@{}}
\toprule
Parameter & Values \\
\midrule
\endhead
\bottomrule
\endlastfoot
Strategies & token\_layer\_group (8 variants), layer\_only (2), token\_only (2) \\
Network size N & 64 \\
Replication R & 2 \\
Churn & 0\%, 5\% \\
Sequence length L & 512 \\
Seeds & 5 (seeds 0--4) \\
\end{longtable}

\textbf{Stage B} (4,050 runs):

\begin{longtable}[]{@{}ll@{}}
\toprule
Parameter & Values \\
\midrule
\endhead
\bottomrule
\endlastfoot
Variants & tb=32 lg=2, tb=32 lg=4, tb=64 lg=2, tb=64 lg=4; layer\_only lb=1 \\
Network size N & 32, 64, 128 \\
Replication R & 1, 2, 3 \\
Churn & 0\%, 5\%, 15\% \\
Sequence length L & 128, 512, 1,024 \\
Seeds & 10 (seeds 0--9) \\
\end{longtable}

\textbf{Stage C} (3,240 runs):

\begin{longtable}[]{@{}ll@{}}
\toprule
Parameter & Values \\
\midrule
\endhead
\bottomrule
\endlastfoot
Variants & tb=32 lg=2, tb=32 lg=4 \\
N, R, Churn, L & Same as Stage B \\
Seeds & 20 (seeds 0--19) \\
\end{longtable}

\section{B. Serialization Format Header Layout}

\begin{verbatim}
Offset  Size   Field
------  ----   -----
0       4 B    Magic: 0x4B564331 ("KVC1")
4       1 B    Version (current: 1)
5       1 B    Compression: 0=none, 1=lz4, 2=zstd
6       1 B    Data type (torch dtype enum)
7       1 B    Reserved
8       4 B    Token start (uint32)
12      4 B    Token end (uint32)
16      4 B    Layer start (uint32)
20      4 B    Layer end (uint32)
24      16 B   Tensor shape: (layers, heads, tokens, head_dim) as 4 x uint32
40      16 B   MD5 checksum of uncompressed payload
56      8 B    Reserved
64      ---    Payload (compressed or raw tensor bytes)
\end{verbatim}

\end{document}
